\scanpage{21}
Donc on a les relations
\[
(13)\quad\left\{\begin{gathered}
\tilde\sigma_0^2=\tilde\sigma_1^2=\tilde\sigma_\infty^2=1\\
\tilde\sigma_0\tilde\sigma_1=\tilde\sigma_1\tilde\sigma_\infty=\tilde\sigma_\infty\tilde\sigma_0
=\overset{\deleted{\sim}}{\rho}\ \deleted{(\unsure{=\rho})}\\
(\overset{\deleted{\sim}}{\rho}{}^3=1\text{ est conséquence des précédentes})
\end{gathered}\right.
\]
\margin{NB
\[
\rho\tilde\sigma_i\rho^{-1}=\tilde\sigma_{\dot\rho i}\quad\text{i.e.}
\]
\[
\begin{aligned}
\rho\tilde\sigma_0\rho^{-1}&=\tilde\sigma_1\\
\rho\tilde\sigma_1\rho^{-1}&=\tilde\sigma_\infty\\
\rho\tilde\sigma_\infty\rho^{-1}&=\tilde\sigma_0
\end{aligned}
\]}
et l'image réciproque $\tilde{\mathcal E}$ de $\varphi(\unsure{\tilde{\mathfrak S}_3|\tilde{\mathcal S}_3})\subset\Gamma$ (ss-groupe d'indice 2) est produit $\frac12$ direct de $\unsure{\tilde{\mathfrak S}_3|\tilde{\mathcal S}_3}$ et de $\pi$. Il faut expliciter les opérations des $\tilde\sigma_i$ sur $\pi$ pour avoir la structure complète du groupe, On trouve que $\tilde\sigma_{\unsure{i|0}}$ est au dessus $\varphi(\tilde\sigma_{\unsure{i|0}})$ i.e. la symétrie p.r. au plan passant par $P,\bar P,\deleted{\unsure{a_i}}i$, qui transforme $l_0$ en $l_0^{-1}$, $l_\infty$ en $l_1^{-1}$, $l_1$ en $l_\infty^{-1}$ :
\margin{Ce sont celles obtenues par transport de structure, via les opérations de $\unsure{\tilde{\mathfrak S}_3|\tilde{\mathcal S}_3}$ sur $(U,P)$.}
\[
(14)\quad\left\{\begin{array}{lll}
\tilde\sigma_0l_0\tilde\sigma_0^{-1}=l_0^{-1}&
\tilde\sigma_0l_1\tilde\sigma_0^{-1}=l_\infty^{-1}&
\tilde\sigma_0l_\infty\tilde\sigma_0^{-1}=l_1^{-1}\\
\tilde\sigma_1l_1\tilde\sigma_1^{-1}=l_1^{-1}&\cdots&\\
\cdots&&
\end{array}\right.
\]
\margin{\uline{Critique} La notation $\tilde\sigma_0$ suggère plutôt que $\tilde\sigma_0$ est au dessus de $\sigma_0$ --- il vaut mieux écrire $\dot\sigma_i'=\dot\sigma_i\tau$ et noter $\tilde\sigma_i'$ au lieu de $\tilde\sigma_i$ $(i\in\{0,1,\infty\})$, ou mieux $\sigma_i'$ (sans $\sim$) --- pour réserver $\tilde\sigma_i'$ à des éléments de $\unsure{S\mathcal E|\mathcal S\mathcal E}$, cf plus bas\ldots}
Les relations (13) (14), plus la relation fondamentale $l_\infty l_1l_0=1$ (1), font une présentation de $\tilde{\mathcal E}$, dont on peut si on veut éliminer $\rho$, en omettant d'écrire la relation qui le définit.

On peut maintenant revenir à nos trois \unsure{miroirs}, \unsure{d.} par
\[
(15)\quad\left\{\begin{aligned}
\tau_0&=\sigma_0\tilde\sigma_0=\tilde\sigma_0\sigma_0\\
\tau_1&=\sigma_1\tilde\sigma_1=\tilde\sigma_1\sigma_1\\
\tau_\infty&=\sigma_\infty\tilde\sigma_\infty=\tilde\sigma_\infty\sigma_\infty
\end{aligned}\right.
\qquad\text{donc}\quad(16)\left\{\begin{aligned}
\tau_0^2&=1\\\tau_1^2&=1\\\tau_\infty^2&=1
\end{aligned}\right.
\]
\mspage{20}
